% Full proof body; requires amsmath and amssymb. Raw H1--H66 retained.
\section{Holonomy descent: the inverse domain, positive root masses, and sharp source realization}
This addendum retains the original \(\chi=\chi_{h,k}\), its quotient
\(E=\mathbb C[S]/(\chi)\), multiplication \(A[P]=[SP]\), the original
source \(g/h\) with \(g=2\xi\), and every measure in (H9)--(H17).
It makes the inverse domain in (H49c) precise, proves the positive-support
and shifted-density assertions used in (H64), and realizes the sharp
quadratic loss with actual compactly supported source columns.
The independent proof review identified the analytic-circle calibration.
The full source-density proof PSD1--41 is an accompanying dependency.

\subsection{The correlation strip and the domain of the quotient inverse}
Fix the original source map \(\Psi:\mathcal P_N\to L^2(\mathbb R;\mathcal K)\),
its Gram \(M=\Psi^*\Psi\succ0\), quotient \(J:\mathcal P_N\twoheadrightarrow E\),
and literal relation injection \(B:\mathbb C^r\hookrightarrow\mathcal P_N\).
Define
\[
 G=(JM^{-1}J^*)^{-1},\quad C=M^{-1}J^*G,\quad M_B=B^*MB.
 \tag{HDI1}
\]
Then \(JC=I\), \(C^*MC=G\), \(B^*MC=0\). For \(r=0\), the
relation metric and its inverse are the unique automorphism of
\(\{0\}\); every correction factoring through \(B\) is the zero
map on its stated domain and codomain. For \(r>0\), \(M_B\succ0\).
Retain \(a>0\), \(M_{a,\pm}=(e^{\pm ar}\Psi)^*(e^{\pm ar}\Psi)\), and
\[
 \kappa_\pm=\sup_{v\ne0}\frac{v^*M_{a,\pm}v}{v^*Mv},\quad
 K_a=\sqrt{\kappa_+\kappa_-},\quad
 \kappa=\sup_{v\ne0}\frac{v^*(M_{a,+}+M_{a,-})v}{v^*Mv}.
 \tag{HDI2}
\]
These retain the factors \(x_k^{\pm2a}\) in the original coordinates.
For \(\mathcal C(v)=\int\Psi(r)^*\Psi(r+v)\,dr\), weighted
Cauchy--Schwarz on two vectors gives
\[
 |\langle x,\mathcal C(v)y\rangle|
 \le K_a e^{-a|v|}\|x\|_M\|y\|_M,\qquad
 \|M^{-1}\mathcal C(nL)\|_M\le K_a e^{-a|n|L}.
 \tag{HDI3}
\]
For \(v>0\), write the integrand as \(e^{-av}\) times the pairing
of \(e^{-ar}\Psi(r)x\) with \(e^{a(r+v)}\Psi(r+v)y\).
Their squared integrals are \(x^*M_{a,-}x\) and \(y^*M_{a,+}y\).
For \(v<0\), use the opposite signs. The dual norm in the retained
\(M\)-metric then proves the operator inequality.
Thus
\[
 \widetilde M(z)=\sum_{n\in\mathbb Z}e^{inz}\mathcal C(nL),\quad
 d_s=\frac{2K_a}{e^{aL-s}-1},\quad
 \|M^{-1}(\widetilde M(z)-M)\|_M\le d_s
 \quad(|\operatorname{Im}z|\le s<aL).
 \tag{HDI4}
\]
The series and its derivative series converge normally on every
closed substrip of \(|\operatorname{Im}z|<aL\): powers of \(|n|\)
are absorbed by the remaining geometric decay. The zero coefficient
is \(M\); summing both nonzero signs proves the displayed bound.
This constructs a holomorphic, periodic extension of the real Gram
coefficients, without conjugating the complex parameter in a Gram.

Choose \(0<s<aL\) with \(d_s<1\). Put
\(T=M^{-1}(\widetilde M-M)\). The maps and their adjoints are
\[
 \begin{array}{lll}
 C:(E,G)\to(\mathcal P_N,M),&
 C^\dagger=G^{-1}C^*M,&C^\dagger C=I,\\
 B:(\mathbb C^r,M_B)\to(\mathcal P_N,M),&
 B^\dagger=M_B^{-1}B^*M,&B^\dagger B=I.
 \end{array}
 \tag{HDI5}
\]
They are isometries onto orthogonal subspaces by (HDI1).
Consequently the four compressions of \(T\) by \(B,C\) have norm
at most \(d_s\), with these domain and target metrics. In particular,
\[
 B^*\widetilde M B=M_B(I+B^\dagger TB),\qquad
 (I+B^\dagger TB)^{-1}=\sum_{j\ge0}(-B^\dagger TB)^j
 \tag{HDI6}
\]
exists and has norm at most \((1-d_s)^{-1}\) on the selected
closed strip. Continuity of \(d_s\), and its strict inequality,
give \(s'>s\) with \(s'<aL\), \(d_{s'}<1\). Hence the inverse is
holomorphic on a neighborhood of the closed strip used by the
contour calculation. The phrase \emph{throughout that strip}
after (H49c) means this selected strip, not the whole larger
correlation strip \(|\operatorname{Im}z|<aL\).

The corresponding Schur complement is
\[
 \widetilde G=C^*\widetilde M C-
 C^*\widetilde M B(B^*\widetilde M B)^{-1}B^*\widetilde M C,
 \qquad
 \|G^{-1}\widetilde G\|_G
 \le1+d_s+\frac{d_s^2}{1-d_s}=\frac1{1-d_s}.
 \tag{HDI7}
\]
The bound follows from (HDI5)--(HDI6), including the two cross
compressions and their actual metrics. On the real axis, subtract
\(B(B^*M_\theta B)^{-1}B^*M_\theta C\) from \(C\). This remains
a section of \(J\) and is \(M_\theta\)-orthogonal to \(\ker J\).
Adding any relation adds precisely its nonnegative squared norm,
so this is the unique minimizing section \(C_\theta\), and its
Gram is the displayed Schur complement.

For Fourier coefficients
\(\widehat G_n=(2\pi)^{-1}\int_0^{2\pi}e^{-in\theta}
\widetilde G(\theta)\,d\theta\), shift to imaginary part \(-s\)
if \(n>0\), and \(+s\) if \(n<0\). The just-proved neighborhood
of holomorphy permits the contour, and periodicity cancels the
two vertical sides. The exponential has modulus \(e^{-s|n|}\).
Therefore
\[
 \|G^{-1}\widehat G_n\|_G\le\frac{e^{-s|n|}}{1-d_s},\quad
 \left\|G^{-1}\left(\frac1m\sum_{j=0}^{m-1}G_{2\pi j/m}
 -\overline G\right)\right\|_G
 \le\frac{2}{(1-d_s)(e^{sm}-1)}.
 \tag{HDI8}
\]
Absolute convergence permits averaging term by term.
The finite character sum keeps precisely the indices divisible
by \(m\); summing their two nonzero tails proves the second bound.

Explicit finite choices exist. Choose \(s_0>0\), \(0<d_0<1\),
\(0<\eta_0<1\), \(0<\varepsilon_0<1-\eta_0^2\), and set
\[
 \begin{split}
 L&=\frac1a\max\left\{s_0+\log(1+2K_a/d_0),
                         \log(1+\kappa/\eta_0)\right\},\\
 m&=\max\left\{1,\left\lceil\frac1{s_0}
 \log\left(1+\frac{2}{(1-d_0)\varepsilon_0}\right)\right\rceil\right\}.
 \end{split}
 \tag{HDI9}
\]
The positive weighted source forms give \(K_a>0\), so \(aL>s_0\).
Monotonicity gives \(d_{s_0}\le d_0\),
\(\delta_L=\kappa/(e^{aL}-1)\le\eta_0\), and error at most
\(\varepsilon_0\) in (HDI8). With the delivered quadratic theorem,
\[
 (1-\eta_0^2-\varepsilon_0)G
 \preceq \frac1m\sum_{j=0}^{m-1}G_{2\pi j/m}
 \preceq(1+\varepsilon_0)G.
 \tag{HDI10}
\]
Its lower factor is positive by the actual selection. These are
finite source-dependent formulas; they do not bound growth of
\(M_{a,\pm}\) with the arithmetic packet. The generator companion
HG gives the combined phase and frequency budget for its operators.

\subsection{Selected-root weights and density on the shifted lattice}
Fix \(L>0\), \(k\ge1\), and \(0\le\theta<2\pi\), and retain
\[
 u_n=\frac{2\pi n+\theta}{L},\quad S_n=k/2+iu_n,\quad
 \nu_n=\frac{2\pi}{L}m_{h,k}(u_n),\quad
 I_\theta=\{n:\nu_n>0\}.
 \tag{HDI11}
\]
The source-density proof PSD1--22 supplies, for each
\(0<b<\pi/2\), the actual finite constants and pointwise bound
\[
 A_b=\|e^{b|\cdot|}w_h\|_\infty,\quad
 B_b=\|e^{b|\cdot|}w_h\|_1,\quad
 m_{h,k}(u)\le A_bB_b^{k-1}e^{-b|u|}.
 \tag{HDI12}
\]
The full PSD source accompanies this addendum; its Gamma factor,
leading coefficient, and cancellation calculation are retained.
The entire nonzero function \(g/h\) has a discrete zero set on
the vertical line, so \(w_h\) is continuous and positive away
from that discrete set. For a fixed \(u\), choose \(t\) outside
the zero sets of both \(w_h(t)\) and \(w_h(u-t)\). An interval
of strictly positive convolution integrand follows by continuity,
so \(m_{h,2}(u)>0\). Induction integrates the preceding everywhere
positive convolution against \(w_h\) on an interval of positivity.
Thus \(I_\theta=\mathbb Z\) for \(k\ge2\).

For \(k=1\), use \(I_\theta\). At a selected root \(\rho\), let
\(d_\rho\) be its full multiplicity in both \(g\) and \(h\). Then
\[
 \begin{split}
 g(s)&=(s-\rho)^{d_\rho}a_\rho(s),\quad
 h(s)=(s-\rho)^{d_\rho}b_\rho(s),\quad
 a_\rho(\rho)b_\rho(\rho)\ne0,\\
 (g/h)(\rho)&=\frac{a_\rho(\rho)}{b_\rho(\rho)}
 =\frac{g^{(d_\rho)}(\rho)}{h^{(d_\rho)}(\rho)}\ne0.
 \end{split}
 \tag{HDI13}
\]
The leading coefficient of \(h\) stays in \(b_\rho\) and its
derivative. In the original monic packet convention
\(\chi_{h,1}=h\). If a nonzero leading coefficient \(a_h\) is
displayed, the monic minimal polynomial is \(a_h^{-1}h\) and its
ideal is \((h)\); the analytic source remains exactly \(g/h\).
Consequently every sampled root \(\rho=S_n\) of \(\chi_{h,1}\)
has the exact positive weight
\[
 \nu_n=\frac1L
 \left|\frac{g^{(d_\rho)}(\rho)}{h^{(d_\rho)}(\rho)}\right|^2>0.
 \tag{HDI14}
\]
This is the full-order cancellation fact that justifies every
root index printed in (H64), even for \(k=1\).

Here is the density proof with the phase retained. Let
\(\mu_n>0\) on \(I\subseteq I_\theta\) have
\(\sum_Ie^{c|u_n|}\mu_n<\infty\) for some \(c>0\), and suppose
\(v\in\ell^2(I,\mu)\) is orthogonal to all polynomial sequences.
For \(0<d<c/2\), Cauchy--Schwarz proves
\[
 \sum_I|v_n|\mu_ne^{d|u_n|}
 \le\|v\|_\mu\left(\sum_Ie^{2d|u_n|}\mu_n\right)^{1/2}<\infty.
 \tag{HDI15}
\]
Therefore \(F(z)=\sum_I\overline v_n\mu_ne^{izu_n}\) is
holomorphic for \(|\operatorname{Im}z|<c/2\), with all
derivative series normally convergent on smaller strips.
The powers are absorbed by the remaining exponential margin.
The original-coordinate polynomial \(((S-k/2)/i)^j\) evaluates
to \(u_n^j\), so every derivative at zero vanishes. The identity
theorem gives \(F=0\). Extend the coefficients by zero off \(I\).
Absolute summability and \(u_n-u_m=2\pi(n-m)/L\) give
\[
 \frac1L\int_0^L F(x)e^{-iu_mx}\,dx
 =\sum_I\overline v_n\mu_n
 \frac1L\int_0^L e^{2\pi i(n-m)x/L}\,dx
 =\begin{cases}\overline v_m\mu_m,&m\in I,\\0,&m\notin I.\end{cases}
 \tag{HDI16}
\]
The left side is zero, so \(v=0\). This proves density.
Apply it to \(\mu=\nu\) on \(I_\theta\), and then to
\(\mu_n=|\chi(S_n)|^2\nu_n\) where \(\chi(S_n)\ne0\).
For both, (HDI12) gives the required exponential moment because
a fixed polynomial is absorbed by the difference of two
positive exponential exponents.

Put \(Z_\theta=\{n\in I_\theta:\chi(S_n)=0\}\) and
\(p_\theta=\prod_{n\in Z_\theta}(S-S_n)\). Coordinate evaluation
is continuous in \(\mathcal H_\theta=\ell^2(I_\theta,\nu)\),
with bound \(|f_n|\le\nu_n^{-1/2}\|f\|_\nu\).
Thus the closure of \(\chi\mathbb C[S]\) vanishes on \(Z_\theta\).
Conversely, if \(f\) vanishes there, division by \(\chi(S_n)\)
on the complement puts \(f/\chi\) in the second weighted space,
with exactly the same squared norm. Polynomial approximation
there, followed by multiplication by \(\chi\), approximates
\(f\) in \(\mathcal H_\theta\). Consequently
\[
 \begin{split}
 \overline{\chi\mathbb C[S]}
 &=\{f\in\mathcal H_\theta:f_n=0\ (n\in Z_\theta)\},\\
 \mathcal T_\theta:=
 \mathcal H_\theta/\overline{\chi\mathbb C[S]}
 &\cong\left(\mathbb C^{Z_\theta},
          \sum_{n\in Z_\theta}\nu_n|a_n|^2\right),\\
 e_\theta:E\twoheadrightarrow\mathcal T_\theta,\quad
 [P]&\longmapsto(P(S_n))_{n\in Z_\theta},\quad
 \ker e_\theta=(p_\theta)/(\chi).
 \end{split}
 \tag{HDI17}
\]
The isometry is restriction, the orthogonal representative of
the Hilbert quotient. The evaluation map is well-defined since
\(\chi\) vanishes at these nodes. The explicit Lagrange polynomial
\(\prod_{m\in Z_\theta,\ m\ne n}(S-S_m)/(S_n-S_m)\)
proves surjectivity. Distinctness of the nodes proves the
divisibility description of the kernel. The empty product is
one and its target is zero. Polynomial cancellation also gives
the exact sequence, equivariant for multiplication by \(S\),
\[
 0\longrightarrow\mathbb C[S]/(\chi/p_\theta)
 \xrightarrow{[Q]\mapsto[p_\theta Q]}E
 \xrightarrow{e_\theta}\mathcal T_\theta\longrightarrow0.
 \tag{HDI18}
\]
No coprimality of \(p_\theta,\chi/p_\theta\) is required.
At a root \(\lambda\) of multiplicity \(d_\lambda\), the original
primary block is \(\mathbb C[\epsilon]/(\epsilon^{d_\lambda})\).
At a sampled root this map is its constant coefficient, with
kernel \(\epsilon\mathbb C[\epsilon]/(\epsilon^{d_\lambda})\);
at an unsampled root the full block is in the kernel. This
follows by substituting \(S=\lambda+\epsilon\) into evaluation,
using the full explicit Chinese-remainder isomorphism in
PSD34--35.

For precision about the original source, let
\(b_n(\theta)=\widehat{\mathcal U_kF_h^{\otimes k}}(u_n)\).
The source-to-sequence dictionary is
\[
 \widehat{I_\theta a}(n)=L^{-1/2}a_nb_n(\theta),\quad
 \|b_n(\theta)\|_{\mathcal K}^2=2\pi m_{h,k}(u_n),\quad
 a_n=\sqrt L\,\frac{\langle b_n(\theta),\widehat f(n)\rangle}
                         {\|b_n(\theta)\|_{\mathcal K}^2}
 \quad(n\in I_\theta).
 \tag{HDI19}
\]
The Fourier basis is
\(L^{-1/2}e^{-i(2\pi n+\theta)r/L}\). Partial Plancherel on the
retained relative variables gives the middle equality, as
proved pointwise in PSD16. Parseval proves the isometry and its
inverse on its range. Polynomial density identifies that range
with the closed polynomial source. On polynomials the map is
exactly
\(I_\theta(P(S_n))=\mathcal Z_{L,\theta}\mathcal U_k\mathcal VP\)
by the plus-sign Fourier convention. All amplitudes and the
factor \(L^{-1/2}\) remain in the comparison.

The measurable field can be specified without an unspecified
choice of fibre coordinates. Embed \(\mathcal T_\theta\) in
\(\ell^2(\mathbb Z)\) by \(a_n\mapsto\sqrt{\nu_n(\theta)}a_n\)
on \(Z_\theta\), zero elsewhere. Its range projection is
\[
 Q_\theta=\operatorname{diag}\!
 \left(1_{\{\nu_n(\theta)>0,\ \chi(S_n(\theta))=0\}}\right)_{n\in\mathbb Z}.
 \tag{HDI20}
\]
Each entry is Borel, since \(\nu_n\) and \(\chi(S_n)\) are
continuous. The sections \(Q_\theta e_n\) form a countable
measurable generating family; approximating any fixed vector
by finitely supported vectors proves strong measurability of
\(Q_\theta v\). A nonzero entry requires a root \(\rho\) with
\(\operatorname{Re}\rho=k/2\) and
\(\theta\equiv L\operatorname{Im}\rho\pmod{2\pi}\).
There are only finitely many such phases. Thus the decomposable
projection \(f(\theta)\mapsto Q_\theta f(\theta)\) on
\(L^2(d\theta/(2\pi);\ell^2(\mathbb Z))\) is zero, because its
squared norm is an integral supported on a finite set.
It follows exactly that
\[
 \int^\oplus\mathcal T_\theta\,\frac{d\theta}{2\pi}=0,\qquad
 E\longrightarrow\int^\oplus\mathcal T_\theta\,\frac{d\theta}{2\pi},
 \quad[P]\mapsto[e_\theta[P]]_{\mathrm{a.e.}},
 \quad\ker=E.
 \tag{HDI21}
\]
The section in this map is Borel by its explicit coordinates
\(1_{Z_\theta}(n)\sqrt{\nu_n(\theta)}P(S_n(\theta))\), and has
zero norm. Its pointwise exceptional maps remain (HDI17)--(HDI18)
with their positive weights; passage to the almost-everywhere
class is the exact arrow removing the measure-zero section.
The accompanying continuous-cochain comparison further records
this operation on its source complex.

On every original support label \(\ell\), each coefficient map
above lifts as \((\ell,v)\mapsto(\ell,fv)\), with
\(\tau\mapsto\tau\). In particular (HDI21) takes a represented
input to the zero in its receiving fibre, with represented
support still present, while external absence remains fixed.
The computed Hilbert kernel does not erase that support.

\subsection{A compactly supported source realizes the full sharp phase loss}
Fix \(L>0\), \(0\le\eta<1\), and the literal mass \(7\).
Choose two nonzero smooth functions \(f_0,g_0\) with disjoint
compact supports in \((0,L)\). Write
\(\alpha=\int|f_0|^2>0\), \(\beta=\int|g_0|^2>0\), and set
\[
 f=f_0/\sqrt\alpha,\quad g_1=g_0/\sqrt\beta,\quad
 \psi_0(r)=\sqrt7\,f(r),\quad
 \psi_1(r)=\sqrt7\left(\eta f(r-L)+\sqrt{1-\eta^2}\,g_1(r)\right).
 \tag{HDI22}
\]
Functions are extended by zero off their supports. The factors
\(\alpha,\beta\) explicitly describe the isometric coordinates
of this constructed source; its Gram retains the mass \(7\).
Disjoint supports give \(\langle f,g_1\rangle=0\) and
\(\langle f,f(\cdot-L)\rangle=0\), while
\(\|f\|=\|g_1\|=\|f(\cdot-L)\|=1\).
Direct integration therefore gives \(M=7I_2\) for
\(\Psi(v_0,v_1)=v_0\psi_0+v_1\psi_1\).
All exponential weighted Grams and derivative moments are
finite, because these are smooth compactly supported columns.

On \(0\le r<L\), only the \(a=0\) terms of \(f,g_1\) and the
\(a=1\) term of \(f(\cdot-L)\) contribute to the actual
holonomy transform \(\sum_a e^{ia\theta}\psi(r+aL)\). Hence
\[
 \mathcal Z_{L,\theta}\psi_0=\sqrt7 f,\quad
 \mathcal Z_{L,\theta}\psi_1=
 \sqrt7(\eta e^{i\theta}f+\sqrt{1-\eta^2}\,g_1),\quad
 M_\theta=7\begin{pmatrix}
 1&\eta e^{i\theta}\\\eta e^{-i\theta}&1
 \end{pmatrix}.
 \tag{HDI23}
\]
Its mean is exactly \(7I_2\). The off-diagonal relative matrix
squares to \(\eta^2I_2\), so its relative eigenvalues are
\(1-\eta,1+\eta\). This is an entire trigonometric matrix
family, obtained from the displayed source with the original
sign convention for the phase.

Take \(J=(1\ \ 0)\), \(B=(0\ \ 1)^T\). Solving \(JC_\theta=1\)
and \(B^*M_\theta C_\theta=0\), and integrating the original
relation discrepancy, gives
\[
 \begin{gathered}
 C=\binom10,\quad C_\theta=\binom{1}{-\eta e^{-i\theta}},
 \quad X_\theta=\eta e^{-i\theta},\\
 G=7,\quad G_\theta=7(1-\eta^2),\quad
 X_\theta^*B^*M_\theta BX_\theta=7\eta^2,\quad
 \mathscr D=7\eta^2.
 \end{gathered}
 \tag{HDI24}
\]
Each formula follows by matrix multiplication, and
\(C-C_\theta=BX_\theta\) is the actual relation identity.
Thus equality in (H33) holds on the full phase circle in this
class of smooth compactly supported Zak sources.
The explicit arrow \(\Psi:\mathbb C^2\to L^2(\mathbb R)\)
and its quotient \(J\) specify this calibration's domain.
They do not identify these columns with a particular arithmetic
packet \(F_h^{\otimes k}\). The arithmetic application continues
to use its original \(\mathcal U_k\mathcal V,J_N,\chi\) and the
theta primitives of (H7), (H26), and (H53)--(H55).
